% Method section for the independent successor study.
% NOT part of iclr_draft.tex. INDEPENDENT_PAPER_PLAN.md forbids replacing the
% predecessor draft's method while its old success tables remain in place.
% Compile standalone, or \input into the successor manuscript.

\section{Method}\label{sec:method}

We adapt the coefficients of a learned dynamics model at the frozen policy's own
action interface. The policy, its weights and its output semantics are untouched.
Everything below acts between the action the policy emits and the action the
robot receives.

\subsection{Interface, state and units}\label{sec:method:interface}

At control step $k$ the frozen policy emits a nominal action $a_{\mathrm{nom}}\in\R^m$
in the interface's own units. The adapter emits $a_{\mathrm{sent}}$, and the robot
executes $a_{\mathrm{sent}}$. We write $\xi$ for the \emph{context}: whatever
measured quantities the basis is allowed to read (configuration, past measured
motion, contact information), which is not in general the estimator state.

Two properties of the interface drive the design, and both were violated by the
predecessor system we build on.

\textbf{The model must carry state.} A correction that is a function of past
\emph{commands} alone cannot represent a plant whose response depends on where the
arm currently is. Measured on the predecessor's own residual, past measured motion
explains $44$--$71\%$ of what its command-only model left unexplained. We therefore
use a descriptor model with an explicit state term.

\textbf{One input map, used twice.} The residual and the allocation must be
computed with the same matrix. When identification and allocation use two
separately fitted input maps, the correction is applied through a different
channel from the one it was estimated in. On our calibration those two maps
disagree by $\|I-M\,\mathrm{pinv}(E)\|_2=2.79$, with $71.6\%$ of the mismatch
off-diagonal, while each is individually well determined.

\subsection{Descriptor model and healthy reference}\label{sec:method:model}

The plant is written in descriptor form,
\begin{equation}
B\,\dot x \;=\; A\,x \;+\; E(\xi,z_E)\,a_{\mathrm{sent}} \;+\; \Phi_d(\xi)\,z_d ,
\label{eq:plant}
\end{equation}
with state $x$, known nominal pair $(A,B)$, input map $E$, and an unknown additive
disturbance expanded in a basis $\Phi_d(\xi)$ with coefficients $z_d$. The input
map carries its own coefficients $z_E$, so loss of actuation authority and additive
disturbance occupy structurally distinct locations: an additive estimate cannot
manufacture a control direction the robot no longer has.

The healthy reference runs the \emph{unmodified} policy action through the
\emph{nominal} model,
\begin{equation}
B\,\dot x_r \;=\; A\,x_r \;+\; E_0\,a_{\mathrm{nom}},
\qquad e \;=\; x - x_r ,
\label{eq:reference}
\end{equation}
so $e$ measures departure from what the frozen policy would have achieved on a
healthy robot, not departure from a hand-specified trajectory. The adapter has no
task-level target of its own.

\subsection{Residual measured from the command actually sent}\label{sec:method:residual}

The prediction residual must subtract the action the robot actually received,
$a_{\mathrm{sent}}$, not the one the policy requested:
\begin{equation}
y \;=\; B\,\dot x - A\,x - E\,a_{\mathrm{sent}} \;\approx\; \Phi(\xi)\,z .
\label{eq:residual}
\end{equation}
Subtracting $a_{\mathrm{nom}}$ instead would feed the adapter's own correction back
in as though it were disturbance, and the loop would chase itself. This is enforced
structurally in our implementation: \code{allocate} and \code{observe} must
alternate, and calling either twice in a row raises rather than silently producing
a residual computed against a stale command.

\subsection{Composite coefficient law}\label{sec:method:law}

The coefficients follow a composite law with a leakage term, a prediction-error
term and a tracking term, with the gain supplied by a Riccati flow:
\begin{align}
\dot z &= -\Lambda z \;+\; P\!\left[\,\Phi^\top R^{-1}\!\left(y-\Phi z\right)
        \;+\; \Phi^\top B^{-\top} L\,e \,\right], \label{eq:zdot}\\
\dot P &= -2\Lambda P \;+\; Q \;-\; P\,\Phi^\top R^{-1}\Phi\,P . \label{eq:pdot}
\end{align}
The same coefficient both predicts the residual and feeds forward into the
correction; the tracking term is what cancels the $e^\top\Phi\tilde z$ cross term
in the Lyapunov ledger and removes the need for persistent excitation. Reading
\eqref{eq:pdot} as a Kalman--Bucy flow gives the tuning dictionary: $\Lambda$ is
mean reversion of the environment, $Q$ how fast it drifts, $R$ representation plus
measurement error. After the deterministic tracking injection $P$ is an adaptation
metric, not a posterior covariance, and we do not report it as one.

\subsection{Allocation}\label{sec:method:allocation}

The correction is the action that cancels the estimated disturbance through the
\emph{same} input map used in \eqref{eq:residual},
\begin{equation}
\hat E\,a_{\mathrm{sent}} \;=\; E_0\,a_{\mathrm{nom}} \;-\; \Phi_d\,\hat z_d .
\label{eq:allocate}
\end{equation}
When the target is not reachable the solve is a constrained least-squares problem
under the interface's box limits, and an infeasible target raises rather than
being silently clipped to something the caller did not ask for. A bounded
correction limit is applied on top; in our experiments it is not the binding
mechanism (Section~\ref{sec:method:status}).

\subsection{Discretisation}\label{sec:method:discrete}

The sampled implementation is derived rather than obtained by replacing
derivatives with differences. We propagate $(\Lambda,Q)$ over one interval by the
Van Loan construction to obtain $(F_z,Q_d)$, apply the measurement update in
Joseph form, and distinguish the window-integrated basis $\Psi$ used for the
prediction residual from the endpoint basis $H_{k+1}$ used for tracking. Using one
where the other belongs changes the fixed point, not merely the transient.

\subsection{Geometry}\label{sec:method:geometry}

Geometry enters as a transformation law on the representation and on every
estimator object, not as an extra network input.

\paragraph{Wrench features.} Contact acts on the arm through generalized forces,
so the disturbance basis is built from a wrench dictionary pulled back through the
contact Jacobians,
\begin{equation}
\Phi_\tau(q) \;=\; \sum_c J_c(q)^\top \Phi_{F,c},
\label{eq:wrench}
\end{equation}
with $J_c$ the world-axis Jacobian at contact $c$ and $\Phi_{F,c}$ in force and
moment units. The output is in generalized-force units and therefore belongs in
torque rows. It cannot be inserted into a first-order position-servo residual
without an explicitly identified conversion; we return to this in
Section~\ref{sec:method:status}.

\paragraph{Equivariance.} For a group element $g$ acting on the context we require
the basis to be equivariant,
\begin{equation}
\Phi(g\cdot\xi) \;=\; \rho_x\,\Phi(\xi)\,\rho_z^\top .
\label{eq:equivariance}
\end{equation}
This is a condition on the \emph{map}, not invariance of a matrix: the family is
equivariant while any individual member may be asymmetric, and a one-sided fault
is carried to its mirror rather than being forbidden. Imposing $\rho_z z=z$ instead
would collapse exactly the coefficients that distinguish the two sides.

Prior specifications of this construction take $G$ to be a finite morphological
group of replicated limbs, which for a single seven-degree-of-freedom arm is
trivial and makes \eqref{eq:equivariance} vacuous. The relevant structure here is
instead a \emph{placement} covariance. Rotate the base mounting and the whole
scene together about the gravity axis, holding the joint coordinates $q$ fixed and
carrying base placement in the context $\xi$. Recomputed world-axis Jacobians then
satisfy $J_c' = D_\theta J_c$ with
$D_\theta=\mathrm{blkdiag}(R_z(\theta),R_z(\theta))$, so \eqref{eq:wrench} obeys
\eqref{eq:equivariance} with
\begin{equation}
\rho_x \;=\; I_{n_q},\qquad \rho_z \;=\; D_\theta ,
\label{eq:yawrep}
\end{equation}
joint torques being frame-independent scalars. Two independent checks confirm it:
on an analytic serial chain the defect of \eqref{eq:equivariance} is
$1.3\times10^{-15}$ while the corresponding \emph{invariance} defect is $5.66$, and
on the actual Panda kinematics over $64$ legal configurations and $8$ yaw angles
the maximum defect is $4.4\times10^{-15}$. The condition is therefore a genuine
constraint rather than one that holds vacuously.

\paragraph{Commutant.} Equivariance of the basis is only consistent if the
estimator and controller transform with it. Every fixed design matrix must lie in
the commutant of the representation acting on its own index,
\begin{equation}
[L,\rho_x]=[R,\rho_x]=0,\qquad [Q,\rho_z]=[\Lambda,\rho_z]=0 ,
\label{eq:commutant}
\end{equation}
so that the adapted response does not depend on the arbitrary orientation of the
world frame. The online estimate itself stays free, transforming as
$\hat z\mapsto\rho_z\hat z$ and $P\mapsto\rho_zP\rho_z^\top$; asymmetric faults
remain representable. For the raw wrench basis $\rho_x=I_{n_q}$ leaves $L$ and $R$
unrestricted, and the constraint falls entirely on the latent gains: under the
gravity-preserving yaw group a symmetric $Q$ drops from $21$ free parameters to
$7$ and a general $\Lambda$ from $36$ to $12$; imposing the full rotation group
would reduce them to $3$ and $4$, and adding reflections to $2$ and $2$, since
parity makes the polar force block and the axial moment block inequivalent. We
enforce \eqref{eq:commutant} at construction, so a violating design matrix raises
rather than being silently projected.

\paragraph{What this does not claim.} \eqref{eq:yawrep} is a covariance of the
basis map over base placements. It is \emph{not} a symmetry of the deployed
experiment, and we claim no nontrivial active symmetry there. Three concrete
facts break one: Panda's first joint is bounded at $\pm2.8973$\,rad, so shifting
it is not a circle action on the legal configuration domain; the interface's
coordinatewise action box is not rotation invariant, failing at quarter-turn
$\pi/4$ on an admissible command; and our own identified nominal model is not
yaw-equivariant, with relative Frobenius defects of $0.080$ and $0.153$ on its
two blocks at $0.37$\,rad. The last of these matters most: imposing
\eqref{eq:commutant} is consistent only alongside a nominal model that respects
the same law, which the current artifact does not. Geometry here is a specified
and algebraically verified construction, and the parameter reductions above are
what it \emph{would} buy once a compatible mechanical model is identified.

\subsection{What is established, and what is not}\label{sec:method:status}

We state the evidence for each component separately, because they are not equally
supported.

\begin{itemize}\itemsep2pt
\item \textbf{Descriptor model, residual convention, composite law, allocation and
discretisation} are implemented and evaluated end to end on a frozen policy under
physical joint faults. Across eight paired cells spanning six joints and two task
suites, the adapter repairs $39$ of $41$ episodes the fault had broken while
breaking $5$ of $119$ that were already succeeding (exact McNemar
$p=1.4\times10^{-7}$). Saturation of the correction limit stays below $10\%$ in
every cell, so the bound is not the mechanism producing the improvement.

\item \textbf{The equivariance and commutant construction} is implemented and
verified algebraically and numerically, as reported above. It has never run in
closed loop, and no nontrivial symmetry of the deployed controller or the frozen
policy is certified.

\item \textbf{The wrench basis \eqref{eq:wrench} is specified but not switched on.}
Our identified model's rows are normalized end-effector motion increments, while
\eqref{eq:wrench} produces generalized forces. The obstruction is not dimensional:
the admissible bridge is a local, single-interval, state-dependent rectangular map
$C_\tau(\xi,a_{\mathrm{sent}})\in\R^{6\times n_q}$ from a held generalized joint
torque to the next six normalized-motion innovation rows, giving the feature
$[\,C_\tau\Phi_\tau\,;\,0\,]$. What is missing is its \emph{identification}. Our
existing rollouts cannot supply it: they carry only persistent single-joint torque
faults, both arms are faulted, and the frozen policy keeps reacting, so a matrix
pooled from them would record a policy and trajectory association rather than a
plant-local mobility. Identifying $C_\tau$ needs a controlled experiment that
holds a known torque over one interval from a matched state, without policy
inference in the loop. Until then the adapter refuses every current identified
model before touching the robot rather than forcing a placement. \emph{All
reported task results use the plain supplied basis; none of them evaluate the
geometric addition.}
\end{itemize}
